\vspace{-2mm}
\subsection{Overall Architecture}
\label{subsec:overall_architecture}
\vspace{-2mm}

For completeness, we discuss other modules in Equiformer here.

\vspace{-2.25mm}
\paragraph{Embedding.}
This module consists of atom embedding and edge-degree embedding.
For the former, we use a linear layer to transform one-hot encoding of atom species.
For the latter, as depicted in the right branch in  Fig.~\ref{fig:equiformer}(c), we first transform a constant one vector into messages encoding local geometry with two linear layers and one intermediate DTP layer and then use sum aggregation to encode degree information~\citep{gin, graphormer_3d}.
The DTP layer has the same form as that in Eq.~\ref{eq:pre_attn}.
We scale the aggregated features by dividing with the squared root of average degrees in training sets so that standard deviation of aggregated features would be close to $1$.
%Including degree information can easily differentiate the two graphs, and therefore, we place edge-degree embedding at the beginning of Equiformer.
%%%The two embeddings are summed to produce final embeddings of input 3D graphs.


\vspace{-2.25mm}
\paragraph{Radial Basis and Radial Function.}
Relative distances $|| \vec{r}_{ij} ||$ parametrize weights in some DTP layers.
%%%To reflect subtle changes in $|| \vec{r}_{ij} ||$, we represent distances with Gaussian radial basis with learnable mean and standard deviation~\citep{schnet} or radial Bessel basis~\citep{dimenet, dimenet_pp}.
To reflect subtle changes in $|| \vec{r}_{ij} ||$, we represent distances with radial basis like Gaussian radial basis~\citep{schnet} and radial Bessel basis~\citep{dimenet, dimenet_pp}.
We transform radial basis with a learnable radial function to generate weights for those DTP layers.
%The function consists of two MLPs with layer normalization and SiLU and a final linear layer.
\revision{The function consists of a two-layer MLP, with each linear layer followed by LN and SiLU, and a final linear layer.}

\vspace{-2.25mm}
\paragraph{Feed Forward Network.}
Similar to Transformers, we use two equivariant linear layers and an intermediate gate activation for the feed forward networks in Equiformer.


\vspace{-2.25mm}
\paragraph{Output Head.}
The last feed forward network transforms features on each node into a scalar.
We perform sum aggregation over all nodes to predict scalar quantities like energy.
Similar to edge-degree embedding, we divide the aggregated scalars with the squared root of average numbers of atoms.


%\begin{enumerate}
%    \item \note{How E(3) features are generated}
%    \item \note{Radial Basis}
%    \item \note{Feed Forward Network}
%    \item \note{Output Head}
%\end{enumerate}